\chapter{Background e stato dell'arte}\label{chap:background}


\noindent Questo capitolo introduce i concetti su cui poggia il lavoro e li colloca rispetto alla
letteratura. Il problema affrontato nasce da una proprietà del dato --- l'incompletezza --- e dalla
possibilità di rimediarvi sfruttando regolarità osservate; il capitolo ricostruisce perciò il
quadro della qualità dei dati e delle sue dimensioni, la natura dei valori mancanti e i meccanismi
che li producono, le famiglie di tecniche di imputazione con le ipotesi che ciascuna richiede, le
dipendenze funzionali rilassate e le misure di similarità che ne definiscono la semantica, la
riduzione di un insieme di regole nella sua accezione classica, e infine la nozione di
\emph{semantic pruning} adottata in questo lavoro, con la sua collocazione rispetto alle linee di
ricerca affini.

% =======================================================================================
\section{Qualità dei dati e sue dimensioni}\label{sec:dq}

La qualità dei dati non è una proprietà intrinseca del dato ma una relazione fra il dato e l'uso
che ne viene fatto: un dataset è di qualità sufficiente rispetto a un compito, e la stessa
raccolta può risultare adeguata per un'analisi descrittiva e inadeguata per una decisione
operativa. La letteratura organizza questa relazione in \emph{dimensioni}, e la classificazione
più diffusa, derivata dallo studio delle aspettative degli utilizzatori, comprende accuratezza,
completezza, consistenza, tempestività e unicità \cite{wang1996beyond}. Le dimensioni non sono
indipendenti --- un dataset completo ma inconsistente è di uso difficile quanto uno accurato ma
incompleto --- e la loro rilevanza dipende dal compito.

\begin{table}[htbp]
\centering
\caption{Dimensioni della qualità dei dati rilevanti per il lavoro. Per ciascuna sono indicati il
difetto corrispondente e il rapporto con il problema affrontato.}
\label{tab:dq-dimensions}
\footnotesize
\begin{tabular}{@{}lp{3.6cm}p{4.6cm}@{}}
\toprule
dimensione & difetto corrispondente & rapporto con questo lavoro \\
\midrule
completezza & valori assenti in celle che dovrebbero essere valorizzate & è la dimensione di
partenza: l'imputazione esiste per rimediare alla sua violazione \\
accuratezza & valore presente ma diverso da quello corretto & è la dimensione che l'imputazione
può migliorare o peggiorare, ed è misurata dal ground truth delle configurazioni remodulate \\
consistenza & violazione di un vincolo fra attributi & è la proprietà che le RFD esprimono in
forma rilassata, e che il sistema usa per rifiutare le imputazioni incoerenti \\
ridondanza & informazione duplicata o non necessaria & è la dimensione su cui agisce la potatura:
un insieme di regole ridondante è un costo senza contropartita informativa \\
\bottomrule
\end{tabular}
\end{table}

Formalmente, la completezza di un dataset $D$ con schema $\mathcal{A} = \{A_1, \dots, A_m\}$ si
misura come la frazione di celle valorizzate sull'insieme delle celle attese,
$\mathrm{comp}(D) = 1 - |\{(t, A_i) : t[A_i] = \bot\}| \,/\, (|D| \cdot m)$; la sua violazione è
un \emph{valore mancante}, che è l'oggetto della sezione seguente. Nel lavoro la completezza è la
dimensione di partenza --- il sistema studiato esiste per rimediare alla sua violazione --- mentre
l'accuratezza è la dimensione lungo la quale si misura se il rimedio migliora o peggiora il dato.
La ridondanza dei metadati, che non compare nelle classificazioni classiche perché riguarda il
processo di arricchimento e non il dato, è il bersaglio della potatura ed è trattata nella
\autoref{sec:rule-reduction}.

Il quadro complessivo è quello dei sistemi di data quality, in cui l'imputazione è una delle
attività di miglioramento e va valutata insieme ai costi che introduce \cite{batini2016data}. È
questa duplice natura --- miglioramento del dato e costo di elaborazione --- a rendere sensato
chiedersi se una parte dell'informazione impiegata sia superflua.

% =======================================================================================
\section{Valori mancanti}\label{sec:missing}

Un valore mancante non è un valore: è l'assenza di un'informazione. La distinzione classica
riguarda il \emph{meccanismo} che lo ha prodotto, poiché da esso dipende la legittimità di
qualunque tecnica di ricostruzione \cite{rubin1976inference}.

\begin{definition}[Meccanismi di mancanza]
Sia $X$ una variabile e $M$ l'indicatore della sua mancanza. Si distinguono:
\begin{itemize}[leftmargin=1.4em]
    \item \textbf{MCAR} (\emph{Missing Completely At Random}): $M$ è indipendente da $X$ e da
    ogni altra variabile osservata;
    \item \textbf{MAR} (\emph{Missing At Random}): $M$ può dipendere dalle variabili osservate,
    ma non dal valore mancante $X$;
    \item \textbf{MNAR} (\emph{Missing Not At Random}): $M$ dipende dal valore mancante stesso.
\end{itemize}
\end{definition}

La distinzione ha una conseguenza pratica diretta: solo sotto MCAR e MAR la distribuzione dei
dati osservati è informativa rispetto al valore mancante, e solo in quei regimi un metodo di
imputazione basato sulle correlazioni osservate può essere corretto in senso statistico
\cite{little2019statistical}. Sotto MNAR la ricostruzione richiede di modellare il meccanismo di
mancanza, e un metodo che non lo faccia produce stime distorte in modo non controllabile dai dati
disponibili.

Nel contesto di questo lavoro la distinzione ha un ruolo operativo preciso. I dataset impiegati
sono \emph{remodulati}: i valori mancanti sono introdotti artificialmente a partire da un dataset
completo, e il valore originale è conservato come ground truth. Il meccanismo è quindi, per
costruzione, noto e controllato. Va però osservato che la rimozione artificiale è progettata per
essere casuale, e ricade pertanto nel caso MCAR: i risultati del lavoro valgono per questo regime,
e non si estendono automaticamente a mancanze strutturali, in cui l'assenza di un valore è essa
stessa un'informazione --- come accade quando un attributo non è applicabile a una parte delle
tuple.

% =======================================================================================
\section{Data imputation}\label{sec:imputation}

L'imputazione consiste nel sostituire un valore mancante con una stima ottenuta dal resto
dell'informazione disponibile. Le famiglie di tecniche si distinguono per la natura
dell'informazione sfruttata e per le ipotesi che devono valere perché la stima sia sensata.

\begin{table}[htbp]
\centering
\caption{Famiglie di tecniche di imputazione. Per ciascuna sono indicati l'informazione
impiegata, l'ipotesi richiesta e il comportamento rispetto alla copertura.}
\label{tab:imputation-families}
\footnotesize
\begin{tabular}{@{}lp{2.6cm}p{3.2cm}p{2.5cm}@{}}
\toprule
famiglia & informazione impiegata & ipotesi & copertura \\
\midrule
statistica & distribuzione marginale o congiunta di una variabile & il modello globale è corretto;
sotto MAR, stima non distorta & totale: un valore viene sempre prodotto \\
apprendimento automatico & relazioni fra attributi apprese dai dati osservati & i dati mancanti
seguono il medesimo processo dei dati osservati & totale \\
similarità & somiglianza fra tuple (\emph{k}-NN) & tuple vicine hanno valori simili & totale o
quasi, secondo la soglia adottata \\
vincoli & dipendenze fra attributi, espresse come regole & il vincolo è valido sul dato reale &
\emph{parziale}: l'imputazione può essere rifiutata \\
\bottomrule
\end{tabular}
\end{table}

I metodi statistici classici --- media, mediana, moda, regressione --- sono semplici e veloci e
assumono una struttura globale che raramente sussiste; l'imputazione multipla per equazioni
concatenate ne rappresenta la versione moderna e condizionata, in cui ogni variabile è stimata a
partire dalle altre in modo iterativo \cite{vanbuuren2011mice}. I metodi di apprendimento
automatico, come le foreste casuali applicate all'imputazione di dati misti, rinunciano a
un'assunzione distributiva e catturano relazioni non lineari, al prezzo di un costo di
addestramento e di una minore interpretabilità \cite{stekhoven2012missforest}. I metodi basati su
similarità copiano il valore da tuple ritenute vicine, con la scelta della metrica e del numero di
vicini come parametri critici \cite{troyanskaya2001missing}. I metodi basati su vincoli
impiegano dipendenze fra attributi --- funzionali, di inclusione, di ordine --- e possono essere
formulati come un problema di riparazione che massimizza la coerenza complessiva
\cite{rekatsinas2017holoclean, fan2012foundations}; è la famiglia a cui appartiene \renuver{}.

La scelta della famiglia non è neutrale rispetto alla valutazione. Un metodo statistico produce
sempre un valore: la sua copertura è per definizione totale, e la sua qualità si misura
sull'accuratezza. Un metodo basato su vincoli può invece \emph{rifiutare} di imputare, quando
nessun candidato soddisfa i vincoli: in quel caso la copertura è una metrica tanto rilevante
quanto l'accuratezza, e le due vanno lette insieme. Questa osservazione è alla base delle metriche
adottate nel \autoref{chap:experiments}, ed è anche la ragione per cui una modifica dell'insieme
di vincoli --- come la potatura studiata in questo lavoro --- può cambiare il risultato in due
direzioni opposte, migliorando la qualità delle imputazioni prodotte e riducendo il numero di
celle che ricevono un valore.

% =======================================================================================
\section{Relaxed Functional Dependencies}\label{sec:rfd}

Una dipendenza funzionale (FD) classica asserisce che due tuple concordanti su un insieme di
attributi $X$ concordano necessariamente su un attributo $A$. Su dati reali l'asserzione è troppo
rigida: il rumore e le eccezioni la falsificano quasi ovunque. Le dipendenze funzionali
\emph{rilassate} introducono forme di tolleranza, tipicamente sulle soglie di confronto e sulle
eccezioni ammesse \cite{caruccio2016rfd}.

\begin{definition}[Relaxed Functional Dependency]\label{def:rfd}
Una RFD è un'espressione della forma
\[
    (A_1 \le \theta_1,\ \dots,\ A_k \le \theta_k) \;\rightarrow\; (A_{\RHS} \le \theta_{\RHS})
\]
dove gli $A_i$ sono attributi del dataset, $\theta_i \ge 0$ sono soglie di tolleranza e
$A_{\RHS}$ è l'attributo di destra. La parte sinistra è detta \LHS{}, la destra \RHS{}.
\end{definition}

La semantica della tolleranza dipende dal tipo dell'attributo: per un attributo numerico
$A_i \le \theta_i$ significa che due tuple possono differire di al più $\theta_i$ su $A_i$; per un
attributo categorico significa che i valori devono coincidere, oppure che la loro distanza di
stringa non supera $\theta_i$. La distinzione è rilevante nel seguito, poiché il sistema studiato
tratta i due casi con procedure diverse.

Due osservazioni sono necessarie per la lettura del lavoro.

\begin{observation}[Ruolo doppio]\label{obs:dual}
Una RFD non è soltanto un descrittore del dataset: è un \emph{oggetto operativo}. Nel sistema
studiato essa partecipa a due procedure distinte, la generazione di candidati per una cella nulla
e la verifica di coerenza di un'imputazione. Le due procedure hanno requisiti opposti rispetto
alla tolleranza: una soglia larga rende la regola permissiva nella generazione e severa nella
verifica. Questo doppio ruolo è il vincolo che determina l'intero problema affrontato
(\autoref{sec:dual-role}).
\end{observation}

\begin{observation}[Soglia nulla]
Il caso $\theta_{\RHS} = 0$ merita attenzione separata. Una RFD con soglia di destra nulla
richiede coincidenza esatta sull'attributo di destra, ed è quella che il sistema tratta con una
scorciatoia dedicata. Nel seguito la distinzione fra regole a soglia nulla e regole a soglia
positiva ricorrerà come variabile esplicativa.
\end{observation}

% =======================================================================================
\section{Ridondanza e riduzione di un insieme di regole}\label{sec:rule-reduction}

L'insieme delle RFD estratte da un dataset è, come ogni insieme di dipendenze, ridondante: alcune
regole sono implicati da altre, e la loro presenza non aggiunge informazione sui vincoli del dato.
La teoria classica delle dipendenze funzionali affronta il problema con la nozione di
\emph{copertura}: due insiemi di dipendenze sono equivalenti se hanno la stessa chiusura, e un
insieme è \emph{minimale} se nessuna delle sue dipendenze è derivabile dalle altre per mezzo degli
assiomi di Armstrong \cite{armstrong1974dependency, maier1983theory}. Il calcolo di una copertura
minimale è un'operazione standard, e la sua applicazione alle RFD è stata studiata come forma di
semplificazione degli insiemi di regole.

Questa letteratura fornisce il termine di paragone naturale del lavoro, ed è importante
chiarire perché non lo risolve. La nozione classica di ridondanza è \emph{dichiarativa}: una
dipendenza è superflua se la sua rimozione non cambia l'insieme dei vincoli soddisfatti. Nel
sistema studiato il ruolo di una regola è invece \emph{operativo}, e dipende da come il programma
la consuma: una regola partecipa alla generazione di un valore e alla verifica di coerenza, e la
scorciatoia per le regole a soglia di destra nulla mostra che due regole equivalenti dal punto di
vista logico possono avere costi ed effetti diversi. Ne segue che la ridondanza classica non
implica la sicurezza della rimozione, e che il problema studiato non è un caso particolare del
calcolo di una copertura minimale.

Il \autoref{chap:theory} formalizza questa distinzione. La conseguenza per la valutazione
sperimentale è che il termine di paragone del lavoro non è una copertura minimale --- che non
garantirebbe l'equivalenza del comportamento --- ma l'\emph{equivalenza osservata} del sistema
prima e dopo la potatura, misurata per confronto degli output. La scelta è più onerosa e più
debole sul piano teorico, ma è l'unica che risponda alla domanda posta: se una regola possa essere
rimossa senza che il sistema se ne accorga.

% =======================================================================================
\section{Misure di similarità e distanza}\label{sec:similarity}

Il confronto fra tuple avviene attributo per attributo, e il risultato complessivo è una
aggregazione dei confronti elementari. Nel sistema studiato l'aggregazione è una media dei
contributi normalizzati sull'ampiezza del lato sinistro della regola. La normalizzazione ha
conseguenze non ovvie: a parità di scostamento assoluto, un contributo è tanto minore quanto più
grande è la soglia di tolleranza dell'attributo, sicché attributi con scale diverse non sono
confrontabili direttamente.

Per gli attributi numerici il confronto elementare è la differenza assoluta, eventualmente
normalizzata sull'intervallo dei valori osservati. Per gli attributi testuali si impiegano misure
di edit distance, fra cui la distanza di Levenshtein, che conta il numero minimo di inserimenti,
cancellazioni e sostituzioni necessari a trasformare una stringa nell'altra
\cite{levenshtein1966binary}; la sua versione normalizzata sulla lunghezza della stringa è
impiegata per rendere confrontabili valori di lunghezza diversa, ed è quella adottata per la
famiglia Physician, i cui attributi testuali designano la stessa entità con forme diverse. Anche
in questo caso la soglia di tolleranza della regola agisce da parametro di accettazione.

La scelta della misura non è neutrale, come mostra il difetto D1 discusso nella
\autoref{sec:d1}: nel sistema studiato l'attributo categorico non riceve la misura prevista, e il
confronto ricade su una distanza di stringa applicata a etichette lunghe e prive di relazione. La
circostanza è rilevante per l'interpretazione dei risultati, perché rende il comportamento del
sistema dipendente da una misura che non è quella dichiarata.

% =======================================================================================
\section{Semantic pruning: definizione operativa e collocazione}\label{sec:semantic-pruning}

Con \emph{semantic pruning} si intende, in questo lavoro, la selezione di un sottoinsieme di un
insieme di metadati --- nel caso in esame, le RFD --- guidata dal \emph{significato} che quei
metadati hanno per il sistema che li consuma, e non da criteri sintattici o statistici sul file.

La distinzione è sostanziale. Un criterio sintattico rimuove regole in base alla loro forma
(duplicati, dimensione del lato sinistro, valore delle soglie); un criterio statistico le rimuove
in base alla loro frequenza o alla loro resa osservata. Un criterio semantico richiede invece di
sapere \emph{che cosa il sistema fa} con ciascuna regola, e quindi di derivare le condizioni di
rimozione dal comportamento effettivo del codice. La qualificazione ``semantico'' non è quindi un
sinonimo di ``euristico'': indica la fonte dell'informazione impiegata, che è la semantica
operativa del consumatore dei metadati e non una proprietà statistica del file.

\begin{definition}[Potatura sicura]\label{def:safe-pruning}
Sia $\mathcal{R}$ un insieme di regole e $\Pi$ il programma che le consuma. Una rimozione
$r \in \mathcal{R}$ è \emph{sicura rispetto a una proprietà $P$} se il comportamento di $\Pi$
rispetto a $P$ è invariato quando $r$ è assente.
\end{definition}

La definizione è parametrica rispetto a $P$, e la scelta di $P$ è il nodo del problema. Nel
sistema studiato si possono isolare almeno due proprietà candidate --- la capacità di generare
candidati e il potere di veto --- e si dimostrerà che esse non sono simultaneamente preservabili
oltre una soglia quantificata (\autoref{sec:negative-theorem}).

Questa impostazione distingue il lavoro da tre alternative presenti in letteratura. La prima è la
potatura \emph{a posteriori}, che misura l'effetto di ogni rimozione eseguendo il sistema:
corretta ma computazionalmente proibitiva, poiché richiede una esecuzione per candidato. La
seconda è la potatura basata su un \emph{proxy} offline, che stima la rilevanza di una regola con
una grandezza calcolata sul dataset: economica, ma --- come si mostrerà nel \autoref{chap:results}
--- potenzialmente fuorviante, al punto da risultare \emph{anti-correlata} con la rilevanza
effettiva. La terza è la selezione guidata dalla semantica \emph{dei nomi} dei metadati, che
impiega modelli linguistici per inferire relazioni fra colonne a partire dalla loro
denominazione \cite{trummer2024llmcorrelations} o per guidare l'esplorazione di fonti dati
eterogenee \cite{freire2024llmdiscovery}: si tratta di una linea di ricerca affine per ambizione
--- impiegare il significato invece della forma --- ma distinta per oggetto, perché riguarda
l'\emph{interpretazione} dei metadati e non la loro \emph{selezione} rispetto al consumo da parte
di un algoritmo noto.

La metodologia adottata cerca una quarta via, intermedia fra le precedenti: derivare dal codice
condizioni \emph{sufficienti}, verificabili offline, la cui validità non dipende da una stima. È
questa la scelta che rende il risultato negativo del \autoref{chap:theory} possibile e
significativo: poiché le condizioni non sono euristiche, il loro limite è un limite del problema e
non dell'implementazione.
